\documentclass{article}

\usepackage{fullpage}
\usepackage{url}
\usepackage{graphicx}
\usepackage{amsmath}
\usepackage{physics}
\usepackage{mathtools}
\usepackage{bbold}

\usepackage{lmodern}
\usepackage[T1]{fontenc}
\usepackage[fontsize=13.5pt]{fontsize}

\DeclareMathOperator{\erfi}{erfi}
\newcommand{\R}{\ensuremath{\mathbb{R}}}
\newcommand{\C}{\ensuremath{\mathbb{C}}}

\title{Quantum Field Theory I Excercise 1}
\author{Idan Stark}
\date{\today}

\begin{document}
\maketitle

\section{The wave functional of the vacuum 1 - canonical quantization}
\subsection{The ground state of the harmonic oscillator}
The Hamiltonian of the QHO
\begin{equation}
    H = \frac{\hat{p}^2}{2M} + \frac{M\omega^2\hat{x}^2}{2}
\end{equation}
Can be rewritten using ladder operators
\begin{equation}
    \hat{a} = \sqrt{\frac{M\omega}{2}}\left(\hat{x} + \frac{i}{M\omega}\hat{p}\right)
    \qquad
    \hat{a}^\dagger = \sqrt{\frac{M\omega}{2}}\left(\hat{x} - \frac{i}{M\omega}\hat{p}\right)
\end{equation}
To get the form:
\begin{equation}
    H = \omega\left(\hat{a}^\dagger\hat{a} + \frac{1}{2}\right)
\end{equation}
This results in the number basis $\ket{n}$ of eigenstates. The lowest eigenstate $\ket{0}$, is the ground state, and is defined by
\begin{equation}
    \hat{a}\ket{0} = 0
\end{equation}
Writing this with the definition of $\hat{a}$ in position space, we get:
\begin{equation}
    \left(x + \frac{1}{M\omega}\partial_x\right)\psi_0(x) = 0
\end{equation}
This is a first order ODE. It is easy to see that it can be solved by a Gaussian:
\begin{equation}
    \psi_0 = Ae^{-x^2/\sigma^2}
\end{equation}
Plugging this into the equation, we get:
\begin{equation}
    \left(x - \frac{2x/\sigma^2}{M\omega}\right)\psi_0(x) = 0
\end{equation}
\begin{equation}
    \sigma = \sqrt{\frac{2}{M\omega}}
\end{equation}
So the ground state is (up to normalization):
\begin{equation}
    \psi_0 = \exp[-\frac{M\omega}{2}x^2]
\end{equation}
\subsection{The free scalar field}
The free scalar field is defined by the action:
\begin{equation}
    S = \int d^4x \left(\frac{1}{2}\partial_\mu\phi \partial^\mu\phi - \frac{1}{2}m^2\phi^2\right)
\end{equation}
Since $\phi$ includes the "position" of the infinite harmonic oscillators, one for every point in space, it is analogous to $\hat{x}$. Then, we can use the Lagrangian definition for momentum and see that $\pi = \frac{\partial \mathcal{L}}{\partial (\partial_0\phi)} = \dot{\phi}$ is analogous to $\hat{p}$. Using this analogy, they must have the same commutation relations for any specific point:
\begin{equation}
    \left[\phi(x), \pi(y)\right] = i\delta(x - y)
\end{equation}
Expanding the kinetic term in the Lagrangian density, we can write the Hamiltonian as:
\begin{equation}
    H = \int d^3x \left[
        \frac{1}{2}\pi^2 + \frac{1}{2}\left(\nabla\phi\right)^2 + \frac{1}{2}m^2\phi^2 
    \right]
\end{equation}
\subsection{The Wave functional of the ground state}
Let us now move to momentum space. The field and the momentum density can be written in term of their Fourier transform $\tilde{\phi}(k)$ and $\tilde{\pi}(k)$:
\begin{equation}
    \phi = \int \frac{d^3k}{(2\pi)^3} \tilde{\phi} e^{ikx}
    \qquad
    \pi = \int \frac{d^3k}{(2\pi)^3} \tilde{\pi} e^{ikx}
\end{equation}
Plugging this into the Hamiltonian gives:
\begin{equation}
    H = \int d^3x \int \frac{d^3k}{(2\pi)^3} \int \frac{d^3p}{(2\pi)^3} \left[
        \frac{1}{2}\tilde{\pi}(k)\tilde{\pi}(p) + \frac{1}{2}(k \cdot p)\tilde{\phi}(k)\tilde{\phi}(p) + \frac{1}{2}m^2\tilde{\phi}(k)\tilde{\phi}(p)
    \right]e^{ix(k + p)}
\end{equation}
The x-integral over the exponent then gives a delta function $(2\pi)^3 \delta(k+p)$. This solves the $p$ integral, giving:
\begin{equation}
    H = \int \frac{d^3k}{(2\pi)^3} \left[
        \frac{1}{2}\tilde{\pi}(k)\tilde{\pi}(-k) + \frac{1}{2}\left(k^2 + m^2\right)\tilde{\phi}(k)\tilde{\phi}(-k)
    \right]
\end{equation}
This is an effective HO in every point in k-space, with effective frequency $\omega_k^2 = k^2 + m^2$. Using this analogy, we can write ladder operators for every $k$ in k-space:
\begin{equation}
    a(k) = \sqrt{\frac{\omega_k}{2}}\left(\tilde{\phi}(k) + \frac{i}{\omega_k}\tilde{\pi}(k)\right)
    \qquad
    a^\dagger(k) = \sqrt{\frac{\omega_k}{2}}\left(\tilde{\phi}(-k) - \frac{i}{\omega_k}\tilde{\pi}(-k)\right)
\end{equation}
And then rewrite the Hamiltonian as:
\begin{equation}
    H = \int \frac{d^3k}{(2\pi)^3} \omega_k \left[a^\dagger(k)a(k) + \frac{1}{2}
    \delta(0)\right]
\end{equation}
Where we used the canonical quantization relations
\begin{equation}
\begin{gathered}
    \left[\tilde{\phi}(k), \tilde{\pi}(p)\right] = \\ =
    \int d^3x e^{-ikx} \int d^3y e^{-iky} \left[\phi(x), \pi(y)\right] = \\ =
    i\int d^3x e^{-ikx} \int d^3y e^{-ipy} \delta(x - y) = i\int d^3x e^{-i(k + p)} = i\delta(k + p)
\end{gathered}
\end{equation}
And so
\begin{equation}
    \left[\tilde{\phi}(k),\tilde{\pi}(-k)\right] = i\delta(0)
\end{equation}
Now we can define the ground state functional $\tilde{\Psi}_0[\tilde{\phi}(k)]$ for the k-th mode, using
\begin{equation}
    a(k)\tilde{\Psi}_0[\tilde{\phi}(k)] = 0
\end{equation}
And writing this using the momentum space definition, using operators $\hat{\phi}(k)$ and $\hat{\pi}(k)$, we get
\begin{equation}
    \left(\hat{\phi}(k) + \frac{i}{\omega}\hat{\pi(k)}\right) \Psi_0[\phi] = 0
\end{equation}
The field operator in just gives the field value:
\begin{equation}
    \hat{\phi}(k) \Psi = \tilde{\phi}(k) \Psi
\end{equation}
We can find $\hat{\pi}$ by imposing the commutation relations
\begin{equation}
    [\hat{\phi}(k), \hat{\pi}(p)] = i\delta(k + p)
\end{equation}
Similar to the HO case, this is solved by
\begin{equation}
    \hat{\pi}(k) = -i\frac{\delta}{\delta \tilde{\phi}(-k)}
\end{equation}
Checking this:
\begin{equation}
\begin{gathered}
    [\hat{\phi}(k), \hat{\pi}(p)]\Psi =
    [\hat{\phi}(k)\hat{\pi}(p) - \hat{\pi}(p)\hat{\phi}(k)]\Psi = \\ =
    -i\tilde{\phi}(k)\frac{\delta \Psi}{\delta \tilde{\phi}(-p)} + i\frac{\delta}{\delta \tilde{\phi}(-p)}(\tilde{\phi}(k)\Psi) = \\ =
    -i\tilde{\phi}(k)\frac{\delta \Psi}{\delta \tilde{\phi}(-p)} + i\tilde{\phi}(k)\frac{\delta \Psi}{\delta \tilde{\phi}(-p)} + i\frac{\delta \tilde{\phi}(k)}{\delta \tilde{\phi}(-p)}\Psi = \\ =
    i\delta(k + p)\Psi
\end{gathered}
\end{equation}
As expected. So, plugging these operators into the ODE for $\Psi_0$, we get:
\begin{equation}
    \left(\tilde{\phi}(k) + \frac{1}{\omega_k}\frac{\delta}{\delta \tilde{\phi}(-k)}\right) \Psi_0 = 0
\end{equation}
We can use $\tilde{\phi}(-k) = \tilde{\phi}^*(k)$ to write
\begin{equation}
    \left|\tilde{\phi}(k)\right|^2 = \tilde{\phi}^*(k)\tilde{\phi}(k) = \tilde{\phi}(-k)\tilde{\phi}(k)
\end{equation}
and get the solution as gaussian $\Psi_0 = Ae^{-\left|\tilde{\phi}\right|^2/\sigma^2}$, similar to the HO:
\begin{equation}
    \tilde{\phi}(k) - \frac{2\tilde{\phi}(k)}{\omega_k\sigma^2} = 0
\end{equation}
\begin{equation}
    \sigma = \sqrt{\frac{2}{\omega_k}}
\end{equation}
Giving us the ground state of the k-th mode:
\begin{equation}
    \tilde{\Psi}_0[\tilde{\phi}(k)] = \exp[-\frac{1}{2}\omega_k\left|\tilde{\phi}(k)\right|^2]
\end{equation}
he total ground state $\Psi_0[\phi(x)]$ is the product of all such modes:
\begin{equation}
    \Psi_0[\phi(x)] = \prod_{k} \exp[-\frac{1}{2}\omega_k\left|\tilde{\phi}(k)\right|^2] = \exp[-\frac{1}{2}\int \frac{d^3k}{(2\pi)^3} \omega_k \left|\tilde{\phi}(k)\right|^2]
\end{equation}
\subsection{The Length Scale}
Using the length scale from QM, only with the newly defined $\hbar\omega_k = \sqrt{m^2c^4 + k^2}$ (where we reinstates the constants $\hbar$ and $c$ for dimensional consistency):
\begin{equation}
    \lambda_k = \sqrt{\frac{\hbar^2}{m\sqrt{m^2c^4 + k^2}}}
\end{equation}
For the $k = 0$ mode, we get the Compton wavelength:
\begin{equation}
    \lambda_c = \frac{\hbar}{mc}
\end{equation}
\section{The wave functional of the vacuum 2 - path integral derivation}
Let us look at the path integral wave functional:
\begin{equation}
    \Psi_{\phi_1}[\phi_2] = \int \mathcal{D} \phi \Big|_{\phi(-T) = \phi_1}^{\phi(0) = \phi_2}e^{iS[\phi]}
\end{equation}
For the free field:
\begin{equation}
    \mathcal{L} = \frac{1}{2}\partial_\mu\phi\partial^\mu\phi - \frac{1}{2}m^2(1 - i\epsilon)^2\phi^2
\end{equation}
Let $\tilde{\phi}$ be the solution to the EOM:
\begin{equation}
    \left[\partial_\mu\partial^\mu + m^2(1 - i\epsilon)\right]\tilde{\phi} = 0
\end{equation}
Moving to k-space\footnote{Since $\tilde{\phi}$ is taken, the Fourier transform of the field is marked using $\tilde{\varphi}$}:
\begin{equation}
    \tilde{\phi}(x) = \int \frac{d^3k}{(2\pi)^3} \tilde{\varphi}(k) e^{ikx}
\end{equation}
We can write the equation for each mode as:
\begin{equation}
    \left[-\partial_0^2 + k^2 + m^2(1 - i\epsilon)\right]\tilde{\varphi} = 0
\end{equation}
This is an ODE solved by the an exponent
\begin{equation}
    \tilde{\varphi} \propto e^{i\omega_kt}
\end{equation}
With the dispersion relation:
\begin{equation}
    \omega_k^2 = k^2 + m^2(1 - i\epsilon)
\end{equation}
The final result is a super position of the two solutions, positive energ (particles) and negative energy (antiparticles):
\begin{equation}
    \tilde{\phi} = \int \frac{d^3k}{(2\pi)^3} \left(
        a_k e^{i\omega_k t} + b_k^\dagger e^{-i\omega_k t}
    \right) e^{ikx}
\end{equation}
Where we named the operator valued prefactors using the same aas the well known ladder operators for particles and antiparticles, and remembered that the field creates antiparticles and annihilated particles.
We can now denote $\phi = \tilde{\phi} + \delta\phi$ and get:
\begin{equation}
\begin{gathered}
    S = \int d^4x \left[
        \frac{1}{2}\partial_\mu\phi\partial^\mu\phi - \frac{1}{2}m^2(1 - i\epsilon)^2\phi^2
    \right] = \\ =
    \frac{1}{2}\int d^4x \left[
        \partial_\mu(\tilde{\phi} + \delta\phi)\partial^\mu(\tilde{\phi} + \delta\phi) - m^2(1 - i\epsilon)^2(\tilde{\phi} + \delta\phi)^2
    \right] = \\ =
    S[\tilde{\phi}] + S[\delta\phi] + \int d^4x \left[
        \partial_\mu\tilde{\phi}\partial^\mu\delta\phi - m^2(1 - i\epsilon)^2\tilde{\phi}\delta\phi
    \right]
\end{gathered}
\end{equation}
The last integral becomes (using integration by parts on the first term and neglecting boundary terms):
\begin{equation}
    \int d^4x \left[
        \partial_\mu\tilde{\phi}\partial^\mu\delta\phi - m^2(1 - i\epsilon)^2\tilde{\phi}\delta\phi
    \right] =
    \int d^4x \left[
        \partial_\mu\partial^\mu\tilde{\phi}\delta\phi - m^2(1 - i\epsilon)^2\tilde{\phi}
    \right]\delta\phi = 0
\end{equation}
Since the term in the brackets is just the EOM. Plugging this into the path integral gives:
\begin{equation}
    \Psi_{\phi_1}[\phi_2] = \int \mathcal{D} \phi \Big|_{\phi(-T) = \phi_1}^{\phi(0) = \phi_2}e^{iS[\tilde{\phi}]}e^{iS[\delta\phi]} = e^{iS[\tilde{\phi}]} \int \mathcal{D} \delta\phi \Big|_{\delta\phi(-T) = 0}^{\delta\phi(0) = 0}e^{iS[\delta\phi]}
\end{equation}
Where we used $\tilde{\phi}(t_1) = \phi_1$ and $\tilde{\phi}(t_2) = \phi_2$. The remaining integral is independent of $\phi_1$ and $\phi_2$ and therefore can be regarded as part of the normalization constant of the probability functional:
\begin{equation}
    \Psi_{\phi_1}[\phi_2] \propto e^{iS[\tilde{\phi}]}
\end{equation}
Writing the action expicitly gives:
\begin{equation}
\begin{gathered}
    S[\tilde{\phi}] = \frac{1}{2}\int_{-T}^0 dt \int d^3x \left[
        \partial_\mu\tilde{\phi}\partial^\mu\tilde{\phi} - m^2(1 - i\epsilon)^2\tilde{\phi}^2
    \right] = \\ =
    S[\tilde{\phi}] = \frac{1}{2}\int_{-T}^0 dt \int d^3x \int \frac{d^3k}{(2\pi)^3}\frac{d^3p}{(2\pi)^3}
    \left[
        k_\mu p^\mu - m^2(1 - i\epsilon)^2
    \right]\tilde{\varphi}(k)\tilde{\varphi}(p)e^{ix(k + p)} = \\ =
    \frac{1}{2}\int_{-T}^0 dt \int \frac{d^3k}{(2\pi)^3}
    \left[
        \omega_k^2 + k^2 + m^2(1 - i\epsilon)^2
    \right]\tilde{\varphi}(k)\tilde{\varphi}(-k) = \\ =
    \int_{-T}^0 dt \int \frac{d^3k}{(2\pi)^3}\omega_k^2
    \left|\tilde{\varphi}(k)\right|^2 = \\ =
    -i\int \frac{d^3k}{(2\pi)^3}\omega_k
    \left(
        \left|\tilde{\varphi}_2(k)\right|^2
        -
        \left|\tilde{\varphi}_1(k)\right|^2
    \right)
\end{gathered}
\end{equation}
Where in the last transition we used the field we wrote in equation (41) to integrate by $t$. The part related to $\phi_1$ is again a constant. As $T \rightarrow -\infty$ that part goes to the vaccum state. We are left with:
\begin{equation}
\begin{gathered}
    \Psi_0[\phi] = \lim_{\epsilon\rightarrow0}\lim_{T\rightarrow-\infty}\exp[i \left(
        -i\int \frac{d^3k}{(2\pi)^3}\omega_k
    \left(
        \left|\tilde{\varphi}(k)\right|^2
        -
        \left|\tilde{\varphi}_1(k)\right|^2
    \right)
    \right)] = \\ =
    \lim_{\epsilon\rightarrow0}\exp[
        \int \frac{d^3k}{(2\pi)^3}\omega_k
    \left(
        \left|\tilde{\varphi}(k)\right|^2
    \right)] = \\ =
    \lim_{\epsilon\rightarrow0}\exp[
        \int \frac{d^3k}{(2\pi)^3}\sqrt{k^2 + m^2(1 + i\epsilon)^2}
    \left(
        \left|\tilde{\varphi}(k)\right|^2
    \right)] = \\ =
    \exp[\int \frac{d^3k}{(2\pi)^3}\sqrt{k^2 + m^2}
    \left(
        \left|\tilde{\varphi}(k)\right|^2
    \right)]
\end{gathered}
\end{equation}
Which is the exact result from the previous question.

\newpage
\section{Gaussian integrals and Wick’s theorem}
\subsection{Basic integral}
Let us define
\begin{equation}
    Z = \int d^nx \exp[-\frac{1}{2}x_iA_{ij}x_j]
\end{equation}
with $A$ being a symmetric strictly positive matrix. This means that $A$ is diagonalizable, so defining $A = P^{-1}DP$, we can write in Dirac notation:
\begin{equation}
    Z = \int d^nx \exp[-\frac{1}{2}\bra{x}P^{-1}DP\ket{x}] = \int d^nx \exp[-\frac{1}{2}\bra{Px}D\ket{Px}]
\end{equation}
So, defining the basis $\ket{y} = \ket{Px}$, and using $\det(P) = 1$ and thus $d^ny = d^nx$, the different degrees of freedom in this basis are seperate, since $D$ is diagonal:
\begin{equation}
\begin{gathered}
    Z = \int d^ny \exp[-\frac{1}{2}\bra{y}D\ket{y}] =
    \int d^ny \exp[-\frac{1}{2}\sum_i D_{ii}y_i^2] = \\ =
    \prod_i \int dy_i \exp[-\frac{1}{2}D_{ii}y_i^2] =
    \prod_i \sqrt{\frac{2\pi}{D_{ii}}} =
    \sqrt{\frac{(2\pi)^n}{\det(D)}} = (2\pi)^{n/2}\left[\det(A)\right]^{-1/2}
\end{gathered}
\end{equation}
Where we used the Gaussain integral and the fact that $\det(A) = \det(P^{-1})\det(D)\det(P) = \det(D)$.

\subsection{Adding term $b$}
Now, let us look integral
\begin{equation}
    Z(b) = \int d^nx \exp[-\frac{1}{2}x_iA_{ij}x_j + b_ix_i]
\end{equation}
for some vector $b$. First let us complete the square:
\begin{equation}
\begin{gathered}
    -\frac{1}{2}x_iA_{ij}x_j + b_ix_i = -\frac{1}{2}\left[x_iA_{ij}x_j - 2 b_ix_i + b_i \Delta_{ij}b_j\right] + \frac{1}{2} b_i \Delta_{ij}b_j = \\ =
    -\frac{1}{2}\left[x_i A_{ij} (x_j - \Delta_{jk} b_k) - b_i (x_i - \Delta_{ij}b_j)\right] + \frac{1}{2} b_i \Delta_{ij}b_j = \\ =
    -\frac{1}{2}\left[(x_i - \Delta_{ik}b_k) A_{ij}(x_j - \Delta_{jl}b_l)\right] + \frac{1}{2} b_i \Delta_{ij}b_j
\end{gathered}
\end{equation}
So, we can define $y_i = x_i - \Delta_{ij}b_j$ and get:
\begin{equation}
    Z(b) = e^{\frac{1}{2} b_i \Delta_{ij}b_j}\int d^ny \exp[-\frac{1}{2}y_iA_{ij}y_j] = (2\pi)^{n/2}\left[\det(A)\right]^{-1/2}e^{\frac{1}{2} b_i \Delta_{ij}b_j}
\end{equation}
\subsection{Correlation functions}
Using the above definition, we can see that:
\begin{equation}
    \langle e^{b_i x_i} \rangle = Z(b)/Z(0) = \exp[\frac{1}{2}b_i\Delta_{ij}b_j]
\end{equation}
Therefore
\begin{equation}
\begin{gathered}
    \langle x_{k_1} x_{k_2} x_{k_3} \cdots  x_{k_n} \rangle = \frac{\partial}{\partial b_{k_1}} \frac{\partial}{\partial b_{k_2}} \frac{\partial}{\partial b_{k_2}} \cdots \frac{\partial}{\partial b_{k_n}} \langle e^{b_i x_i} \rangle \Big|_{b=0} = \\ =
    \frac{\partial}{\partial b_{k_1}} \frac{\partial}{\partial b_{k_2}} \frac{\partial}{\partial b_{k_2}} \cdots \frac{\partial}{\partial b_{k_n}} \exp[\frac{1}{2}b_i\Delta_{ij}b_j] \Big|_{b=0}
\end{gathered}
\end{equation}
Applying the derivatives one by one, we get terms that are the exponent multiplied by some power of $b$. At the ends, all terms where that power is not zero vanish, when setting $b = 0$. This means the only terms that survive are the terms where half the derivatives acted on the exponent, adding the prefactor $\Delta_{ik_n}b_i$ and the other half derived the prefactors, leaving the term with $\Delta_{k_nk_m}$. The assignment $b = 0$ at the end then causes the exponent to be $1$, which means the total result is:
\begin{equation}
    \langle x_{k_1} x_{k_2} x_{k_3} \cdots  x_{k_n} \rangle = \sum_{all \ possible \ pairings \ P} \Delta_{k_{P_1}k_{P_2}} \cdots \Delta_{k_{P_{n-1}}k_{P_n}}
\end{equation}

\section{Gaussian Integration of Grassmann Numbers}
\subsection{One Complex Grassmann variable}
Using the Taylor expansion of Grassmann exponent
\begin{equation}
    e^{\theta} = 1 + \theta
\end{equation}
We can write the Basic Gaussian as
\begin{equation}
    \int d\theta^* d\theta \theta \theta^*e^{-\theta^* a \theta} = \int d\theta^* d\theta \theta \theta^*(1 - \theta^* a \theta) = \int d\theta^* d\theta \theta \theta^*(1 + a \theta \theta^*)
\end{equation}
Since every Grassmann number squared is zero, we get:
\begin{equation}
    \int d\theta^* d\theta \theta \theta^*(1 + a \theta \theta^*) =
    \int d\theta^* d\theta \theta \theta^* + a \left(\theta \theta^*\right)^2 = 
    \int d\theta^* d\theta \theta \theta^* 
\end{equation}
Finally, using
\begin{equation}
    \int d\theta \theta = 1
\end{equation}
We can seperate the variables to get:
\begin{equation}
    \int d\theta^* d\theta \theta \theta^* =
    \int d\theta^* \left(\int d\theta \theta\right) \theta^* = \int d\theta^* \theta^* = 1
\end{equation}
\subsection{Many Complex Grassmann variables}
Now, let us look at the integral with $N$ complex Grassmann variables
\begin{equation}
    \prod_{i=1}^{N} \left(\int d\theta_i^* d\theta_i\right)\theta_j\theta_k^* \exp[-\sum_{l,m=1}^N \theta_l^*A_{lm}\theta_m]
\end{equation}
Using the same logic as in the c-numbers case, we can diagonalize $A_ij$ to get $A = U^{-1}DU$, then define $\eta_i = U_{ij}\theta_j$, seperating the variables:
\begin{equation}
\begin{gathered}
    U^{-1}_{jm}U_{nk}\prod_{i=1}^{N} \left(\int d\eta_i^* d\eta_i\right)\eta_m\eta_n^* 
    \exp[-\sum_{i=1}^N \eta_l^*D_{ll}\eta_l]
\end{gathered}
\end{equation}
The above integral is seperable, leaving two cases. In the $m = n$ case, for all $i \ne n$, the integral is just a simple Gaussian giving $D_{ii}$, while the $m$ integral gives the integral from the last section, resulting in $1$.
\newline
In the $m \ne n$ case, for all $i \ne m,n$, the integral is again just a simple Gaussian, giving $D_{ii}$, while the $n$ and $m$ integrals give:
\begin{equation}
    \int d\eta_m^*d\eta_m \eta_m[1 + D_{mm}\eta^*_m\eta_m] = 0
\end{equation}
Since int he first term, the integral over $\eta^*_m$ is zero and on the second term there is a $\eta_m^2$ which is also zero. This is of course true for both indices for both $\eta$ and $\eta^*$. And so, the result is that in the sum over $n$ only $n = m$ survives, leaving
\begin{equation}
        \sum_m U^{-1}_{jm}U_{mk}\frac{\prod_i D_{ii}}{D_{mm}} = \det A\sum_m\frac{U^{-1}_{jm}U_{mk}}{D_{mm}}
\end{equation}
Now, let us notice that:
\begin{equation}
    \sum_m\frac{U^{-1}_{jm}U_{mk}}{D_{mm}} = \sum_m U^{-1}_{jm}D_{mm}^{-1}U_{mk} = A^{-1}_{jk}
\end{equation}
And so we get:
\begin{equation}
    \det(A)A^{-1}_{jk}
\end{equation}
\end{document}